\documentclass[12 pt, a4paper]{article}% тип документа, размер шрифта
\usepackage{cmap}	
\usepackage{hyperref}
\hypersetup{
    colorlinks=true,
    linkcolor=blue,
    urlcolor=blue,
}
\usepackage{mathtext}
\usepackage[T2A]{fontenc}%поддержка кириллицы в ЛаТеХ
\usepackage[utf8]{inputenc}%кодировка
\usepackage[english,russian]{babel}
\usepackage{indentfirst}

\usepackage{amsmath,amsfonts,amssymb,amsthm,mathtools} % AMS
\usepackage{amsmath}%удобная вёрстка многострочных формул, масштабирующийся текст в формулах, формулы в рамках и др.
\usepackage{amsfonts}%поддержка ажурного и готического шрифтов — например, для записи символа {\displaystyle \mathbb {R} } \mathbb {R} 
\usepackage{amssymb}%amsfonts + несколько сотен дополнительных математических символов
\frenchspacing%запрет длинного пробела после точки
\usepackage{setspace}%возможность установки межстрочного интервала
\usepackage{indentfirst}%пакет позволяет делать в первом абзаце после заголовка абзацный отступ
\onehalfspacing%установка полуторного интервала по умолчанию
\usepackage{graphicx}%подключение рисунков
\graphicspath{{images/}}%путь ко всем рисункам
\usepackage{caption}
\usepackage{float}%плавающие картинки
\usepackage{tikz} % это для чудо-миллиметровки
\usepackage{pgfplots}%для построения графиков
\pgfplotsset{compat=newest, y label style={rotate=-90},  width=10 cm}%версия пакета построения графиков, ширина графиков
\usepackage{pgfplotstable}%простое рисование табличек
\usepackage{lastpage}%пакет нумерации страниц
\usepackage{comment}%возможность вставлять большие комменты
\usepackage{python}%питончик
\usepackage{float}
%%%%% ПОЛЯ
\setlength\parindent{0pt} 
\usepackage[top = 2 cm, bottom = 2 cm, left = 1.5 cm, right = 1.5 cm]{geometry}
\setlength\parindent{0pt}
%%%%% КОЛОНТИТУЛЫ
\usepackage[shortlabels]{enumitem}

\usepackage{fancybox, fancyhdr}
\pagestyle{fancy} 
\fancyhead[L]{Легких Никита Б05-301}
\fancyhead[C]{2.4.1}
% ЛЮБАЯ ДОПОЛНИТЕЛЬНАЯ ИНФОРМАЦИЯ
\fancyfoot[R]{ФПМИ МФТИ}
\renewcommand{\footrulewidth}{0.3 mm}
 

\usepackage{graphicx}
\graphicspath{{pictures/}}
\DeclareGraphicsExtensions{.pdf,.png,.jpg}
\usepackage{wrapfig}
%%% Заголовок
%%% Новые команды
\newcommand{\z}[2]{{{\vspace{0.6cm} \textbf{{Задача {#1}}} \textit{{#2}\\ }}} }
\newcommand{\otv}{{\vspace{0.3cm} \textbf{Ответ: }}}
\newcommand{\uk}{\underline{\textit{Указание.}} }
\newcommand{\prim}{\underline{\textit{Примечание.}} }
\newcommand{\sol}[1]{{{\vspace{0.3cm} \textbf{{Задача {#1}} }\\ }}}



\begin{document} % конец преамбулы, начало документа

\begin{center}
\section*{Работа 2.4.1}
    \section*{Определение теплоты испарения жидкости}
\end{center}
\textbf{Цель работы:} 1) измерение давления насыщенного пара жидкости при разной температуре; 2) вычисление по полученным данным теплоты испарения с помощью уравнения Клапейрона-Клаузиуса. \\ \\
\textbf{В работе используются:} термостат, герметический сосуд, заполненный водой, отсчётный микроскоп.

\section*{Теоретическая справка}
Испарением называется переход вещества из жидкого в газообразное состояние. При испарении с поверхности вылетают молекулы, которые обладают достаточной кинетической энергией для совершения работы против атмосферы и преодоления молекулярных сил сцепления. Вследствие вылета высокоэнергетических молекул жидкость остывает и чтобы это предовратить мы подводим тепло с помощью калориметра. \\
Для вычисления интересной нам теплоты парообразования воды будем пользоваться косвенным методом, так как прямой нам недоступен из-за сложности контроля теплопотерь. Наш метод будет основан на формуле Клапейрона–Клаузиуса: 
\begin{equation}
    \frac{dP}{dT} = \frac{L}{T (V_2 - V_1)} 
\end{equation}
Здесь $P$ — давление насыщенного пара жидкости при температуре $T$ , $T$ — абсолютная температура жидкости и пара, $L$ — теплота испарения жидкости, $V_2$ — объем пара, $V_1$ — объем жидкости. Найдя из опыта $\frac{dP}{dT}$, $T$, $V_2$ и $V_1$, можно определить $L$ путем расчета. При этом величины $L$, $V_2$ и $V_1$ в формуле (1) должны относиться к одному и тому же количеству вещества, поэтому для простоты будем относить их к одному молю. \\
Так как в нашем приборе измерения производятся при давлениях ниже атмосферного, поэтому задача упрощается. В таблице ниже представлены вычесленные величины для различных жидкостей: 
\begin{figure}[h]
\centering
\includegraphics[width=0.6\linewidth]{2.4.1 Определение теплоты испарения жидкости/табличка вычисленных величин.png}
\label{fig:mpr}
\end{figure}

Отсюда видно, что отношение $V_1/V_2 \leq 0.05$ для исследуемой жидкости, поэтому мы в праве заменить $(V_2 - V_1)$ на $V_2$, и для простоты заменим $V_2$ на просто $V$. \newpage
Теперь обратимся к уравнению Ван-дер-Ваальса: 
\begin{equation}
    (P+\frac{a}{V^2})(V-b) = RT
\end{equation}
Из таблицы видно что $b \approx V_1 \Rightarrow b$ можно пренебречь по сравнению с $V$, в пренебрежение членом $a/V^2$ по сравнению с $P$ вносит ошибку менее 3\%. Значит мы в праве пользоваться уравнением идеального газа:
\begin{equation}
    V = \frac{RT}{P}
\end{equation}
Подставляя (3) в (1), пренебрегая $V_1$ и разрешая уравнение относительно $L$, найдем
\begin{equation}
    L = \frac{mRT^2}{\mu mP} \frac{dP}{dT} = - \frac{R}{\mu} \frac{d(\ln{P})}{d(1/T)}
\end{equation}
С помощью нашей установки и приведения линеаризации найдем нужное отношение $d(\ln{P})/d(1/T)$ \\
\section*{Установка}
\begin{figure}[h]
\centering
\includegraphics[width=0.3\linewidth]{2.4.1 Определение теплоты испарения жидкости/установка1.png}
\label{fig:mpr}
\end{figure}
Схема установки изображена на рисунке 1. Наполненный водой резервуар 1 играет
роль термостата. Нагревание термостата производится спиралью 2, подогреваемой электрическим током. Для охлаждения воды в термостате через змеевик 3 пропускается водопроводная вода. Вода в термостате перемешивается воздухом, поступающим через трубку 4. Температура воды измеряется термометром 5. В термостат погружен запаянный прибор 6 с исследуемой жидкостью. Над ней находится насыщенный пар (перед заполнением прибора воздух из него был откачан). Давление насыщенного пара определяется по ртутному манометру
соединенному с исследуемым объемом. Отсчет показаний манометра производится при помощи микроскопа. На рисунке 2 приведена более полная схема такой же установки, но с использованием современного термостата. Установка включает термостат A, экспериментальный прибор B и отсчетный микроскоп C. Экспериментальный прибор B представляет собой емкость 12, заполненную водой. В нее погружен запаянный прибор 13 с исследуемой жидкостью 14. Перед заполнением исследуемой жидкости воздух из запаянного прибора был удален, так что над жидкостью находится только её насыщенный пар. Давление пара определяется по ртутному манометру 15, соединенному с емкостью 13. Численная величина давления измеряется по разности показаний отсчетного микроскопа 16, настраиваемого последовательно на нижний и верхний уровни столбика ртути манометра. Показания микроскопа снимаются по шкале 17.
\begin{figure}[h]
\centering
\includegraphics[width=0.7\linewidth]{2.4.1 Определение теплоты испарения жидкости/установка2.png}
\label{fig:mpr}
\end{figure}
Главная проблема нашей установки в том что мы измеряет не ту температуру, хоть и близкую к ней, чтобы уменьшить эту ошибку нагревание будет делать очень медленным (2-4 минуты на 1$^oC$) и будем промерять и в обратную сторону (с шагом в 2$^oC$)
\section*{Задание}
\begin{enumerate}
    \item $\Delta H_0 \approx 1$ мм
    \item 
    \begin{center}
        \begin{tabular}{|l|l|l|l|l|l|l|l|l|l|l|}
        \hline
        $T, ^oC$ & 21.06 & 22.07 & 23.05 & 24.09 & 25.08 & 26.09 & 27.03 & 28.03 & 29.08 & 30.06 \\ 
        \hline
        $H,$ см & n/a   & n/a   & n/a   & 1.5   & 1.95  & 2.10  & 2.20  & 2.38  & 2.66  & 2.76  \\ 
        \hline
        \end{tabular}
    \end{center}
    \begin{center}
    \begin{tabular}{|l|l|l|l|l|l|l|l|l|l|l|}
    \hline
    $T, ^oC$ & 31.06 & 32.06 & 33.04 & 34.07 & 35.08 & 36.08 & 37.03 & 38.06 & 39.01 & 40.05 \\ \hline
    $H,$ см & 2.94  & 3.14  & 3.39  & 3.59  & 3.88  & 4.05  & 4.27  & 4.57  & 4.79  & 5.00     \\ \hline
    \end{tabular}
    \end{center} 
    \item 
    \begin{center}
    \begin{tabular}{|l|l|l|l|l|l|l|l|l|l|l|l|}
    \hline
    $T, ^oC$ & 40.05 & 38.02 & 36.00 & 34.05 & 32.01 & 30.11 & 28.02 & 26.17 & 23.70 & 21.80 & 20.00 \\ \hline
    $H$, см & 5.00  & 4.60  & 4.12  & 3.75  & 3.34  & 2.95  & 2.63  & 2.32  & 2.00  & 1.69  & 1.53  \\ \hline
    \end{tabular}
    \end{center}
    \item Заметим что: $\ln{P} = \ln{\rho g} + \ln{H}$, а нас интересует лишь угол наклона графика, поэтому нам достаточно построить график в координатах: $\ln{H}$ от $1/T$ и он получится один в один что $\ln{P}$ от $1/T$:\newpage
    \begin{figure}[h]
    \centering
    \includegraphics[width=0.7\linewidth]{2.4.1 Определение теплоты испарения жидкости/графичек.png}
    \label{fig:mpr}
    \end{figure}
    
    Из графика выше: $$k_{нагрева} = -2.19\pm0.07 \;K\cdot10^3$$
    $$k_{охлаждения} = -2.31\pm0.13 \;K\cdot10^3$$
    Значения погрешностей оценивались через формулу МНК: $\Delta k = \sqrt{\frac{<\ln{H}^2>-<\ln{H}>^2}{<1/T^2>-<1/T>^2} - k^2}/\sqrt{N}$
    Пересчитываем значения $k$ в $L \rightarrow$ $$k_{нагрева} \rightarrow L_{нагрева} = -\frac{R}{\mu}k_{нагрева} = 1.012 \; МДж/кг$$ $$\Delta k_{нагрева} \rightarrow \Delta L_{нагрева} = (-\frac{R}{\mu}k_{нагрева})\cdot \epsilon_{k_{нагрева}} = 0.03 \; МДж/кг$$ 
    $$k_{охлаждения} \rightarrow L_{охлаждения} = -\frac{R}{\mu}k_{охлаждения} = 1.067 \; МДж/кг$$ $$\Delta k_{охлаждения} \rightarrow \Delta L_{охлаждения} = (-\frac{R}{\mu}k_{охлаждения})\cdot \epsilon_{k_{охлаждения}} = 0.06 \; МДж/кг$$
\end{enumerate}
\section*{Вывод}
Мы получили значения: $L_{охлаждения} = 1.067\pm0.06 \; МДж/кг$ и $L_{нагрева} = 1.012\pm0.03 \; МДж/кг$, видно что они очень близки друг к другу: $\frac{L_{охлаждения}}{L_{нагрева}} = 1.054$. Но к сожалению оба значения далеки от верного для нашей жидкости, которая судя по справочнику равна: $2.1\pm0.2 \; МДж/кг$, ошибки в обратке данных не обнаружено, значит проблема в методике наших измерений.

\end{document}